\documentclass[11pt, a4paper]{article}

% Gemeinsame Style-Definitionen (TEXINPUTS muss templates/ enthalten — siehe scripts/build_tex.py)
% bfw-style.tex — gemeinsame Style-Definitionen für alle Handouts/Aufgabenhefte
% Voraussetzung: Kompilation mit xelatex (wegen fontspec)

\usepackage[ngerman]{babel}
% Babel-ngerman macht " zu einem aktiven Shorthand (z.B. für "a -> ä). In unseren
% Texten verwenden wir " als normales Zeichen — daher Shorthand komplett ausschalten.
\shorthandoff{"}
\usepackage{geometry}
\geometry{a4paper, top=2.4cm, bottom=2.4cm, left=2.3cm, right=2.3cm,
          headheight=15pt, headsep=0.7cm, footskip=1.1cm}

\usepackage{fontspec}
% Fallback-Kette: Source Sans Pro -> Calibri -> Segoe UI -> Latin Modern Sans
% (Latin Modern ist immer mit MiKTeX vorhanden)
\IfFontExistsTF{Source Sans Pro}{%
  \setmainfont{Source Sans Pro}[Ligatures=TeX]%
}{\IfFontExistsTF{Calibri}{%
  \setmainfont{Calibri}[Ligatures=TeX]%
}{\IfFontExistsTF{Segoe UI}{%
  \setmainfont{Segoe UI}[Ligatures=TeX]%
}{%
  \setmainfont{Latin Modern Sans}[Ligatures=TeX]%
}}}
\IfFontExistsTF{Source Code Pro}{%
  \setmonofont{Source Code Pro}[Scale=0.92]%
}{\IfFontExistsTF{Consolas}{%
  \setmonofont{Consolas}[Scale=0.92]%
}{%
  \setmonofont{Latin Modern Mono}[Scale=0.92]%
}}

% Gesperrte Versalien für Labels/Titelbalken (Kapitälchen-Ersatz, funktioniert
% mit jeder OpenType-Schrift — Latin Modern Sans hat keine echten Kapitälchen)
\newcommand{\bfwlabelfont}{\footnotesize\bfseries\addfontfeatures{LetterSpace=8.0}}

\usepackage{xcolor}
% Farbpalette: hell, freundlich, modern
\definecolor{bfwPrimary}{HTML}{0EA5A4}    % Petrol/Türkis - Primär
\definecolor{bfwPrimaryDark}{HTML}{0F766E} % dunkler Petrol für Headings
\definecolor{bfwAccent}{HTML}{F59E0B}     % Amber - Akzent
\definecolor{bfwSuccess}{HTML}{10B981}    % Grün - Erfolg / Lösung
\definecolor{bfwWarn}{HTML}{F97316}       % Orange - Achtung
\definecolor{bfwInk}{HTML}{0F172A}        % fast Schwarz
\definecolor{bfwInkSoft}{HTML}{475569}    % gedämpftes Grau
\definecolor{bfwBgSoft}{HTML}{F8FAFC}     % off-white
\definecolor{bfwBgTeal}{HTML}{ECFEFF}     % helles Türkis
\definecolor{bfwBgAmber}{HTML}{FEF3C7}    % helles Gelb
\definecolor{bfwBgRose}{HTML}{FFE4E6}     % helles Rosé für Eselsbrücken
\definecolor{bfwTabHead}{HTML}{E4F3F2}    % zarte Kopfzeilen-Hinterlegung für Tabellen

\usepackage[colorlinks=true, linkcolor=bfwPrimaryDark, urlcolor=bfwPrimaryDark]{hyperref}
\usepackage{lastpage}
\usepackage{enumitem}
\setlist[itemize]{leftmargin=*, itemsep=2pt, topsep=2pt}
\setlist[enumerate]{leftmargin=*, itemsep=2pt, topsep=2pt}

\usepackage[most]{tcolorbox}
\usepackage{booktabs}
\usepackage{colortbl}
\usepackage{tikz}
\usetikzlibrary{decorations.pathmorphing, positioning, shapes.geometric, arrows.meta}

% ---------------------------------------------------------------------------
% Einheitliches Tabellen-Design
%   - etwas mehr Zeilenluft, dezente graue Linien (wirkt auch mit \hline ruhiger)
%   - Kopfzeile einer Tabelle bei Bedarf zart hinterlegen: \rowcolor{bfwTabHead}
% ---------------------------------------------------------------------------
\renewcommand{\arraystretch}{1.25}
\arrayrulecolor{bfwInkSoft!50}
\setlength{\heavyrulewidth}{1pt}
\setlength{\lightrulewidth}{0.5pt}

% ---------------------------------------------------------------------------
% Boxen — klare farbige Titelbalken mit gesperrten Versal-Labels
% (Titel bleibt per Option überschreibbar: \begin{merksatz}[title={...}])
% ---------------------------------------------------------------------------
\tcbset{bfwbox/.style={%
  enhanced, breakable,
  boxrule=0.5pt, arc=2.5pt,
  left=10pt, right=10pt, top=8pt, bottom=8pt,
  toptitle=4pt, bottomtitle=4pt,
  titlerule=0pt,
  fonttitle=\bfwlabelfont,
  coltitle=white
}}

% Merkkästchen — wichtige Definition / Kernpunkt
\newtcolorbox{merksatz}[1][]{%
  bfwbox,
  colback=bfwBgTeal,
  colframe=bfwPrimaryDark,
  colbacktitle=bfwPrimaryDark,
  title={MERKE},
  #1
}

% Eselsbrücke — Mnemoniks
\newtcolorbox{eselsbruecke}[1][]{%
  bfwbox,
  colback=bfwBgRose,
  colframe=bfwAccent!85!black,
  colbacktitle=bfwAccent!85!black,
  title={ESELSBRÜCKE},
  #1
}

% Beispiel-Box
\newtcolorbox{beispiel}[1][]{%
  bfwbox,
  colback=bfwBgSoft,
  colframe=bfwInkSoft,
  colbacktitle=bfwInkSoft,
  title={BEISPIEL},
  #1
}

% Lösungs-Box (für Aufgabenhefte)
\newtcolorbox{loesung}[1][]{%
  bfwbox,
  colback=bfwSuccess!7,
  colframe=bfwSuccess!65!black,
  colbacktitle=bfwSuccess!65!black,
  title={LÖSUNG},
  #1
}

% Achtung-Box — typische Fehler / Stolperfallen
\newtcolorbox{achtung}[1][]{%
  bfwbox,
  colback=bfwWarn!6,
  colframe=bfwWarn!85!black,
  colbacktitle=bfwWarn!85!black,
  title={ACHTUNG},
  #1
}

% Formel-Box — für zentrale Formeln (groß, ruhig, ohne Titelbalken)
\newtcolorbox{formelbox}[1][]{%
  enhanced,
  colback=bfwBgTeal,
  colframe=bfwPrimaryDark,
  boxrule=1pt, arc=3pt,
  left=10pt, right=10pt, top=6pt, bottom=6pt,
  #1
}

% Glossar-Eintrag
\newcommand{\glossareintrag}[2]{%
  \par\noindent\textbf{\color{bfwPrimaryDark}#1} \enspace #2\par\smallskip
}

% ---------------------------------------------------------------------------
% Abschnitts-Design: farbiges Nummern-Quadrat + feine Linie unter \section,
% dezent farbige \subsection/\subsubsection
% ---------------------------------------------------------------------------
\usepackage{titlesec}
\newcommand{\bfwSecBadge}[1]{%
  \begin{tikzpicture}[baseline={([yshift=-0.62em]n.center)}]
    \node[fill=bfwPrimaryDark, text=white, font=\normalsize\bfseries,
          minimum width=1.55em, minimum height=1.55em, inner sep=1pt,
          rounded corners=1.5pt] (n) {#1};
  \end{tikzpicture}}
\titleformat{\section}
  {\Large\bfseries\color{bfwPrimaryDark}}
  {\bfwSecBadge{\thesection}}{0.6em}{}
  [\vspace{-0.2em}{\color{bfwPrimary!60}\titlerule[0.8pt]}]
\titleformat{\subsection}
  {\large\bfseries\color{bfwPrimaryDark}}
  {\textcolor{bfwPrimary}{\thesubsection}}{0.5em}{}
\titleformat{\subsubsection}
  {\normalsize\bfseries\color{bfwInk}}
  {\textcolor{bfwPrimary}{\thesubsubsection}}{0.4em}{}
\titlespacing*{\section}{0pt}{1.6em}{1em}
\titlespacing*{\subsection}{0pt}{1.2em}{0.5em}

% TikZ-Stil "skizzenhaft" — etwas weicher als Rough.js, aber stabil
\tikzset{
  roughbox/.style = {
    draw=bfwPrimaryDark, thick, fill=bfwBgTeal,
    rounded corners=4pt, inner sep=6pt
  },
  roughaccent/.style = {
    draw=bfwAccent, thick, fill=bfwBgAmber,
    rounded corners=4pt, inner sep=6pt
  },
  roughline/.style = {
    draw=bfwInkSoft, thick, -{Stealth[length=2.5mm]}
  }
}

% Bullet-Symbole (bewusst kein FontAwesome — vermeidet Font-Installations-Probleme)
\newcommand{\faLightbulb}{$\bullet$}
\newcommand{\faBookmark}{$\bullet$}

% ---------------------------------------------------------------------------
% Kopf-/Fußzeile
%   Kopf links:  Wortlogo „Wirtschaftsfachwirt"
%   Kopf rechts: Dokument-Kurztext, gesetzt via \kopfzeile{...}
%   Fuß links:   © Can Siebert · 2026 — Fuß rechts: Seite X von Y
%   Titelseite (Seite 1) bleibt ohne Kopf-/Fußzeile (\handouttitel setzt
%   \thispagestyle{empty}).
% ---------------------------------------------------------------------------
\usepackage{fancyhdr}
\newcommand{\bfwKopfzeileText}{}
\newcommand{\kopfzeile}[1]{\renewcommand{\bfwKopfzeileText}{#1}}
\pagestyle{fancy}
\fancyhf{}
\fancyhead[L]{\small\bfseries\color{bfwPrimaryDark}Wirtschaftsfachwirt}
\fancyhead[R]{\small\color{bfwInkSoft}\bfwKopfzeileText}
\renewcommand{\headrulewidth}{0.6pt}
\renewcommand{\headrule}{{\color{bfwPrimary}\hrule width\headwidth height 0.6pt}}
\fancyfoot[L]{\small\color{bfwInkSoft}\copyright\ Can Siebert\,\textperiodcentered\,2026}
\fancyfoot[R]{\small\color{bfwInkSoft}Seite \thepage\ von \pageref*{LastPage}}
% Auch \thispagestyle{plain} (z.B. durch \maketitle) soll leer bleiben:
\fancypagestyle{plain}{\fancyhf{}\renewcommand{\headrulewidth}{0pt}\renewcommand{\headrule}{}}

% ---------------------------------------------------------------------------
% \handouttitel{Obertitel}{Untertitel}{Kurs-Label}{Termin-Zeile}
%   Einheitlicher professioneller Titelblock: Dachzeile, farbige Headline,
%   Untertitel, farbige Trennlinie und dezente Meta-Zeile (Kurs/Termin/Dozent).
%   Ersetzt \title/\maketitle. Direkt nach \begin{document} aufrufen.
% ---------------------------------------------------------------------------
\newcommand{\handouttitel}[4]{%
  \thispagestyle{empty}%
  \begingroup
  \setlength{\parindent}{0pt}%
  {\bfwlabelfont\color{bfwPrimary}WIRTSCHAFTSFACHWIRT (IHK)\par}
  \vspace{0.55em}
  {\huge\bfseries\color{bfwPrimaryDark}#1\par}
  \vspace{0.5em}
  {\large\color{bfwInkSoft}#2\par}
  \vspace{0.9em}
  {\color{bfwPrimary}\rule{\linewidth}{2pt}}\par
  \vspace{0.55em}
  {\small\color{bfwInkSoft}#3\ \,\textperiodcentered\,\ #4\ \,\textperiodcentered\,\ Dozent: Can Siebert\par}
  \endgroup
  \vspace{1.4em}%
}


\title{\vspace*{-1.5cm}\Huge\bfseries\color{bfwPrimaryDark}Lektion~2: Wettbewerbspolitik, Staatseingriffe \& Volkswirtschaftliche Gesamtrechnung\\[0.4em]
  \large\color{bfwInkSoft}\textnormal{Vom funktionierenden Wettbewerb zum gesamtwirtschaftlichen Gleichgewicht}}
\author{\textsf{Wirtschaftsfachwirt --- 1.1 Volkswirtschaftliche Grundlagen \textperiodcentered{} \small\copyright\ Can Siebert\,\textperiodcentered\,2026}}
\date{}

\begin{document}
\maketitle
\thispagestyle{empty}

\vspace{1em}
\begin{merksatz}[title={\bfseries Worum geht es in dieser Lektion?}]
Basis der sozialen Marktwirtschaft ist ein funktionierender Markt mit ausreichendem
Wettbewerb. Da ein sich selbst überlassener Markt die Tendenz zeigt, Wettbewerb zu
reduzieren, greift der Staat regulierend ein: durch \emph{Wettbewerbspolitik} (GWB,
Bundeskartellamt) und durch \emph{Eingriffe in die Preisbildung} (Höchst- und
Mindestpreise, Subventionen, Steuern). Ob diese Politik erfolgreich ist, misst die
\emph{Volkswirtschaftliche Gesamtrechnung} (BIP, BNE, Volkseinkommen) --- und die Ziele,
an denen sich die Wirtschaftspolitik ausrichten muss, formuliert das Stabilitätsgesetz
im \emph{magischen Viereck}. Alle vier Themen sind Standardstoff der IHK-Prüfung.
\end{merksatz}

\tableofcontents
\vspace{1em}
\hrule
\vspace{1em}

\section{Wettbewerb und Wettbewerbspolitik}

\subsection{Was ist Wettbewerb?}

Unter \textbf{Wettbewerb} versteht man allgemein die \emph{Rivalität der Marktteilnehmer
unter- und gegeneinander}, um ihre Ziele zulasten der Konkurrenten durchzusetzen. Der
Wettbewerb sichert u.\,a. die persönliche und wirtschaftliche Freiheit, die
Güterverteilung und Bedarfsdeckung über den Marktmechanismus und den ökonomischen
Einsatz der Produktionsfaktoren.

\subsection{Funktionen des Wettbewerbs}

Idealtypisch erfüllt der Wettbewerb folgende Funktionen:

\begin{itemize}
  \item \textbf{Freiheitsfunktion:} Jeder kann seine Fähigkeiten und Mittel für frei
    gewählte Zwecke einsetzen (Freiheit der Unternehmertätigkeit, Freiheit der
    Konsumwahl) --- freier Wettbewerb schafft mehr Handlungsalternativen.
  \item \textbf{Kontrollfunktion:} Anbieter und Nachfrager kontrollieren sich gegenseitig;
    jeder muss sich an der Leistung der Konkurrenz messen lassen. Die wirtschaftliche
    Macht des Einzelnen wird begrenzt.
  \item \textbf{Steuerungsfunktion (Lenkungsfunktion):} Mit zwei Unterfunktionen ---
    \emph{Koordinations- und Anpassungsfunktion} (Wettbewerb beschleunigt die Koordination
    von Angebot und Nachfrage; langfristig wird nur angeboten, was der Nachfrage
    entspricht) und \emph{Allokationsfunktion} (die Produktionsfaktoren Arbeit, Boden,
    Kapital, Know-how werden so kombiniert, dass ein Maximum an Produktivität entsteht).
  \item \textbf{Innovationsfunktion:} Wettbewerb schafft Anreize, Produkte zu verbessern,
    die Produktpalette zu erneuern, kostengünstiger zu produzieren und den technischen
    Fortschritt rasch zu nutzen.
  \item \textbf{Auslesefunktion:} Auslese der Marktteilnehmer nach wirtschaftlicher
    Leistungsfähigkeit --- stets kommen die leistungsfähigsten Anbieter und Nachfrager
    zum Zug.
  \item \textbf{Verteilungsfunktion:} Die Entlohnung der Produktionsfaktoren entspricht
    deren Leistung --- leistungsorientierte Einkommensverteilung.
\end{itemize}

\begin{eselsbruecke}
\textbf{F-K-S-I-A-V} --- „\textbf{F}reier \textbf{K}onkurrenzkampf \textbf{s}ichert
\textbf{i}mmer \textbf{a}lle \textbf{V}orteile``: \textbf{F}reiheit, \textbf{K}ontrolle,
\textbf{S}teuerung, \textbf{I}nnovation, \textbf{A}uslese, \textbf{V}erteilung ---
die sechs Funktionen des Wettbewerbs.
\end{eselsbruecke}

\subsection{Ziele und Instrumente der Wettbewerbspolitik}

Weil unvollkommene Märkte zu Wettbewerbsverzerrungen führen und Märkte die Tendenz haben,
Wettbewerb abzubauen, steuert der Staat im Rahmen der \textbf{Ordnungspolitik} gegen.
Schwerpunkte: freier Zugang zu Märkten, Abbau von Handelshemmnissen, Verhinderung
wettbewerbsbeschränkender Verhaltensweisen. Leitbild ist heute nicht mehr die
vollständige Konkurrenz, sondern der Erhalt von \emph{Qualität und Intensität} des
Wettbewerbs --- Wissen, Kreativität und Innovationen gelten als treibende Marktkräfte.

Die beiden gesetzlichen Grundlagen in Deutschland:

\begin{itemize}
  \item \textbf{UWG} --- Gesetz gegen den unlauteren Wettbewerb: garantiert fairen
    Wettbewerb und faire Marktpraktiken (Schutz der \emph{Art und Weise} des Wettbewerbs).
  \item \textbf{GWB} --- Gesetz gegen Wettbewerbsbeschränkungen („Kartellgesetz``):
    unterbindet Kartelle, Absprachen und Preisbindungen und wirkt der Konzentration der
    Wirtschaft entgegen (Schutz des \emph{Bestands} von Wettbewerb).
\end{itemize}

Über die Einhaltung des GWB wacht das \textbf{Bundeskartellamt}. In der EU liegt die
Wettbewerbspolitik bei der \textbf{EU-Kommission} (Kommissar für Wettbewerb);
Rechtsgrundlagen sind der \textbf{AEUV} (Vertrag über die Arbeitsweise der EU) und die
europäische \textbf{Fusionskontrollverordnung (FKVO)}. Durch die GWB-Novellen (7.~Novelle
2005: Angleichung an EU-Recht; 8.~Novelle 2013: Neustrukturierung der Missbrauchsaufsicht,
Angleichung der Fusionskontrolle; 9.~Novelle 2017: Umsetzung der
EU-Kartellschadenersatzrichtlinie) ähneln sich deutsches und europäisches Recht sehr.
\emph{EU-Recht steht über nationalem Recht} --- es greift, wenn Wettbewerbsbeschränkungen
den Handel zwischen Mitgliedstaaten beeinträchtigen. Ein weltweites Wettbewerbsrecht oder
Kartellamt existiert nicht.

\begin{merksatz}[title={\bfseries Die drei Säulen des GWB}]
\textbf{1. Kartellverbot} (\S\,1 GWB) \quad \textbf{2. Missbrauchsaufsicht}
(\S\S\,19\,ff. GWB) \quad \textbf{3. Fusionskontrolle} (\S\,35 GWB).
Merkhilfe: \textbf{Ka-Mi-Fu}.
\end{merksatz}

\subsection{Das Kartellverbot (\S\,1 GWB)}

\textbf{Kartelle} sind Kooperationen von Unternehmen auf Basis (schriftlicher oder
mündlicher) vertraglicher Vereinbarungen, um die Wettbewerbsposition zu verbessern. Die
Beteiligten bleiben \emph{rechtlich selbstständig}, schränken ihre \emph{wirtschaftliche}
Selbstständigkeit aber teilweise ein. Nach \S\,1 GWB sind Vereinbarungen zwischen
Unternehmen, Beschlüsse von Unternehmensvereinigungen und aufeinander abgestimmte
Verhaltensweisen \textbf{verboten}, wenn sie eine \emph{Verhinderung, Einschränkung oder
Verfälschung des Wettbewerbs} bewirken. Das europarechtliche Pendant ist
\textbf{Art.~101 AEUV}. Verbotene Kartellarten insbesondere:

\begin{center}
\begin{tabular}{p{4.2cm} p{9.8cm}}
\toprule
\textbf{Kartellart} & \textbf{Inhalt der Absprache}\\
\midrule
Preiskartell & Absprache einer einheitlichen Preispolitik\\
Submissionskartell & Vereinbarung von Preisuntergrenzen bzw. Festlegung, welches Kartellmitglied bei öffentlichen Ausschreibungen zum Zug kommt\\
Gebietskartell & Aufteilung des Absatzmarktes\\
Quotenkartell & künstliche Verknappung des Angebots durch Absprache von Produktionsquoten\\
\bottomrule
\end{tabular}
\end{center}

\subsubsection{Ausnahmen vom Kartellverbot}

Obwohl Kartelle grundsätzlich verboten sind, gibt es \textbf{Ausnahmen} in drei Gruppen:

\begin{enumerate}
  \item \textbf{Besondere Ausnahmeregelungen nach dem GWB:} z.\,B.
    \emph{Mittelstandskartelle} (\S\,3 GWB) --- Vereinbarungen zwischen kleinen und
    mittleren Unternehmen, die wegen ihres begrenzten Potenzials Wettbewerbsnachteile
    gegenüber Großunternehmen haben.
  \item \textbf{EU-Gruppenfreistellungsverordnungen:} Bestimmte Gruppen von Vereinbarungen
    sind freigestellt, weil sie sich trotz gewisser Wettbewerbsbeschränkung positiv auf
    den Markt auswirken; Einzelfallprüfung entfällt.
  \item \textbf{Einzelfreistellungen (Legalausnahmen, \S\,2 GWB / Art.~101 Abs.~3 AEUV):}
    Vereinbarungen sind freigestellt, wenn sie (a) den Verbraucher angemessen am Gewinn
    beteiligen, (b) der Verbesserung der Warenerzeugung oder -verteilung dienen oder
    (c) zur Förderung des technischen Fortschritts beitragen --- \emph{ohne} dass der
    Wettbewerb zwischen den Beteiligten ausgeschaltet wird.
\end{enumerate}

Nach dem \textbf{Selbstprüfungssystem} müssen die Unternehmen selbst einschätzen, ob ihr
Verhalten mit dem Kartellrecht vereinbar ist --- Anmelde- und Überprüfungsverfahren beim
Kartellamt entfallen (weniger Bürokratie, mehr Eigenverantwortung). Bei Verstößen drohen
empfindliche Strafen (Bußgelder, Gewinnabschöpfung).

\subsection{Missbrauchsaufsicht (\S\S\,19\,ff. GWB)}

Trotz Wettbewerbspolitik entstehen immer mehr Unternehmen mit erheblicher Marktmacht.
Über \emph{marktbeherrschende} Unternehmen übt das Bundeskartellamt die
\textbf{Missbrauchsaufsicht} aus. Zwei Missbrauchsformen:

\begin{itemize}
  \item \textbf{Ausbeutungsmissbrauch:} Forderung überhöhter Preise von Kunden oder
    extremer Preisnachlässe von abhängigen Zulieferern (viele Zulieferer sind so stark an
    einen Kunden gebunden, dass ihre Existenz bedroht ist, wenn sie ihn verlieren).
  \item \textbf{Behinderungsmissbrauch:} Marktmacht wird gezielt eingesetzt, um aktuelle
    oder potenzielle Konkurrenten zu behindern, zu verdrängen oder ihnen den Marktzugang
    zu versperren --- z.\,B. über \emph{Kampfpreise (Dumping)} oder die Verweigerung des
    Zugangs zu Infrastruktureinrichtungen.
\end{itemize}

\subsection{Fusionskontrolle (\S\,35 GWB)}

Während die Missbrauchsaufsicht \emph{bestehende} marktbeherrschende Unternehmen betrifft,
ist die \textbf{Fusionskontrolle} auch eine \emph{vorbeugende} Maßnahme: Geplante Fusionen
müssen dem Bundeskartellamt \emph{angezeigt} werden. Es kann den Zusammenschluss
\textbf{untersagen}, wenn die Gefahr besteht, dass eine marktbeherrschende Stellung
entsteht. Der Bundeswirtschaftsminister kann einen untersagten Zusammenschluss dennoch
erlauben (\textbf{Ministererlaubnis}, \S\,42 GWB), wenn er von gesamtwirtschaftlichem
Vorteil ist und die marktwirtschaftliche Ordnung nicht gefährdet wird.

\begin{beispiel}
Drei Bauunternehmen sprechen vor einer öffentlichen Ausschreibung ab, wer das günstigste
Angebot abgibt und den Auftrag erhält $\Rightarrow$ verbotenes
\textbf{Submissionskartell} (\S\,1 GWB). Ein marktbeherrschender Onlinehändler verkauft
gezielt unter Einstandspreis, um einen neuen Konkurrenten zu verdrängen $\Rightarrow$
\textbf{Behinderungsmissbrauch} (Dumping). Zwei große Lebensmittelketten wollen
fusionieren und hätten danach regional über 50~\% Marktanteil $\Rightarrow$ Anzeigepflicht,
das Bundeskartellamt kann die Fusion \textbf{untersagen}.
\end{beispiel}

\section{Eingriffe des Staates in die Preisbildung}

In einer Marktwirtschaft bildet sich der Preis durch Angebot und Nachfrage. In der
\emph{sozialen} Marktwirtschaft übernimmt der Staat regulierende Aufgaben, um
Benachteiligungen von Anbietern oder Nachfragern auszugleichen. Grundsätzlich gilt:

\begin{merksatz}
\textbf{Marktkonforme (indirekte) Eingriffe} beeinflussen Angebot und Nachfrage, ohne den
Marktpreismechanismus außer Kraft zu setzen (z.\,B. Steuern, Zölle, Subventionen).
\textbf{Marktkonträre (direkte) Eingriffe} greifen unmittelbar in die Preisbildung ein
und hebeln den Preismechanismus aus (z.\,B. Höchst- und Mindestpreise).
\end{merksatz}

\subsection{Indirekte (marktkonforme) Maßnahmen}

\begin{itemize}
  \item \textbf{Verbrauchssteuern:} Höhere Verbrauchssteuern (z.\,B. Tabaksteuer)
    verteuern das Produkt $\Rightarrow$ höherer Gleichgewichtspreis, geringere
    Gleichgewichtsmenge. So kann Verbraucherverhalten aus umwelt- und
    gesundheitspolitischen Gründen gelenkt werden (weniger Raucher). Mit der geringeren
    Menge sinken allerdings auch die Steuereinnahmen. Steuersenkungen wirken umgekehrt:
    niedrigerer Gleichgewichtspreis, höhere Menge.
  \item \textbf{Zölle:} Senkung von Einfuhrzöllen erhöht die Importe (wirkt gegen
    Preissteigerungen im Inland); Erhöhung von Einfuhrzöllen verteuert Importwaren und
    schützt die einheimische Industrie. Die Gesetzgebungs- und Ertragskompetenz für Zölle
    liegt bei der \emph{EU}.
  \item \textbf{Subventionen und Steuervergünstigungen (Anbieterförderung):} Staatliche
    Finanzhilfen ermöglichen es Anbietern, zu einem Preis unterhalb der
    Produktionskosten zu verkaufen --- z.\,B. zur Markteinführung neuer Produkte
    (Windkraftanlagen). \emph{Problematisch bei Dauersubvention:} Sie behindert sämtliche
    Preisfunktionen; der künstlich niedrige Preis sendet falsche Signale, das Ausscheiden
    kostenungünstiger Anbieter wird verhindert, Produktionsfaktoren bleiben in nicht
    (mehr) gefragten Produkten gebunden.
  \item \textbf{Transferleistungen (Subvention der Nachfrager):} Der Staat gibt Haushalten
    zusätzliche Mittel, damit sie ein bestimmtes Gut erwerben können --- z.\,B.
    \emph{Wohngeld} oder \emph{Kindergeld}. Die Finanzierung von Subventionen und
    Transfers tragen die Steuerzahler insgesamt.
  \item \textbf{Mengenmaßnahmen:} Regulierung der Angebots- und Nachfragemenge, z.\,B.
    Einfuhrkontingente, Einfuhrverbote, staatliche Vorratshaltung.
\end{itemize}

\subsection{Direkte (marktkonträre) Maßnahmen: Mindest- und Höchstpreise}

\subsubsection{Mindestpreis}

Ein \textbf{Mindestpreis} ist eine staatlich verordnete \emph{Preisuntergrenze}, die
nicht unterschritten werden darf; er liegt \emph{über} dem Gleichgewichtspreis.
\textbf{Ziel: Schutz der Anbieter.} Folgen: Das Angebot wird ausgeweitet, die Nachfrage
sinkt $\Rightarrow$ \textbf{Angebotsüberschuss} ($x_N < x_0 < x_A$). Die Überschüsse
stellen den Staat vor neue Probleme --- er muss mengenregulierend eingreifen
(Kontingente, Prämien für Angebotsreduzierung, Flächenstilllegung). Bei
Mengenkontingenten kann ein \textbf{grauer Markt} entstehen, auf dem die Mehrproduktion
unter dem festgesetzten Mindestpreis angeboten wird.

\begin{beispiel}
\textbf{EU-Agrarmarkt:} Landwirten wurden jahrzehntelang Mindestpreise mit
Abnahmegarantien zugestanden --- Folge: massive Überproduktion („Butterberge``,
„Milchseen``). Der Staat musste die Überschüsse durch Einlagerung, Verarbeitung,
Veräußerung oder Vernichtung beseitigen (erhebliche Kosten) und das Angebot durch
Produktionsquoten und Stilllegungsprämien eindämmen.
\textbf{Gesetzlicher Mindestlohn:} Er hebelt die Findung des Gleichgewichtslohns aus und
kann zu einem dauerhaften Angebotsüberhang an Arbeitskräften (Arbeitslosigkeit) führen,
wenn Arbeitgeber Arbeitskräfte durch Technik substituieren oder Arbeitsumfänge ins
Ausland verlagern.
\end{beispiel}

\subsubsection{Höchstpreis}

Ein \textbf{Höchstpreis} ist eine staatlich verordnete \emph{Preisobergrenze}; er liegt
\emph{unter} dem Gleichgewichtspreis. \textbf{Ziel: Schutz der Nachfrager.} Beispiele:
Mietpreisbindung im sozialen Wohnungsbau, \emph{Mietpreisbremse} auf angespannten
Wohnungsmärkten. Folgen: Die Nachfrage steigt, das Angebot sinkt $\Rightarrow$
\textbf{Nachfrageüberschuss} ($x_A < x_0 < x_N$). Damit der Höchstpreis wirkt, sind
weitere Maßnahmen nötig: \emph{Subventionen}, damit überhaupt ein Angebot vorhanden ist,
sowie \emph{Kontingentierung und Rationierung} der knappen Menge. Häufig bilden sich
\textbf{Schwarzmärkte}, auf denen die Güter zu höheren Preisen verkauft werden. Selten
kann das verminderte Angebot über \emph{Importe} gelöst werden, wenn ausländische
Anbieter aufgrund von Kostenvorteilen unter dem Höchstpreis anbieten können.

\begin{eselsbruecke}
\textbf{MI}ndestpreis liegt \textbf{ü}ber, \textbf{HÖ}chstpreis liegt \textbf{unter} dem
Gleichgewichtspreis --- sonst wären sie wirkungslos! Mindestpreis schützt
\textbf{A}nbieter (beide mit Vokal-Anfang: \emph{Mi--A}), Höchstpreis schützt
\textbf{N}achfrager. Mindestpreis $\Rightarrow$ \textbf{A}ngebotsüberschuss,
Höchstpreis $\Rightarrow$ \textbf{N}achfrageüberschuss.
\end{eselsbruecke}

\begin{merksatz}[title={\bfseries Beurteilung staatlicher Preiseingriffe (prüfungsrelevant!)}]
Staatliche Preisfestsetzungen setzen die Funktionen des Marktpreises (Signal, Ausgleich,
Allokation, Selektion) außer Kraft: Es entstehen dauerhafte Überschüsse bzw. Knappheiten,
Folgekosten für den Staat, Fehlallokation von Produktionsfaktoren sowie graue und schwarze
Märkte. Marktkonforme Eingriffe (Steuern, Subventionen, Transfers) sind daher
ordnungspolitisch vorzuziehen --- aber auch Dauersubventionen verzerren die Preissignale.
\end{merksatz}

\section{Volkswirtschaftliche Gesamtrechnung (VGR)}

\subsection{Aufgaben der VGR}

Die \textbf{VGR} ist das statistische Instrument zur zahlenmäßigen Erfassung
gesamtwirtschaftlicher Vorgänge. Sie gibt \emph{rückschauend} für eine abgeschlossene
Periode Auskunft über \textbf{Entstehung, Verwendung und Verteilung} des
Bruttoinlandsprodukts (BIP) und des Bruttonationaleinkommens (BNE). Grundlage ist die
\textbf{Kreislauftheorie} (Wirtschaftskreislauf) mit den Sektoren private Haushalte,
Unternehmen, Staat, Kapitalsammelstellen und Ausland. Die Daten dienen u.\,a. zur
Ermittlung von Wirtschaftswachstum und Strukturwandel, als Konjunkturbarometer, als
Grundlage wirtschaftspolitischer Entscheidungen, als Anhaltspunkt für Lohnverhandlungen
der Tarifpartner und für den internationalen Vergleich. Für EU-Staaten gilt das
\textbf{Europäische System volkswirtschaftlicher Gesamtrechnungen (ESVG 2010)} --- es
sichert die Vergleichbarkeit der Daten zwischen den EU-Ländern.

\subsection{BIP und BNE: Inlandskonzept vs. Inländerkonzept}

\begin{itemize}
  \item Das \textbf{Bruttoinlandsprodukt (BIP)} misst den Wert der Güter und
    Dienstleistungen, die innerhalb eines Zeitraums \emph{innerhalb der geografischen
    Grenzen} einer Volkswirtschaft produziert wurden (\textbf{Inlandskonzept}) ---
    bewertet zu Marktpreisen. Es spielt keine Rolle, ob die Produzenten ihren Wohnsitz
    im Inland haben.
  \item Das \textbf{Bruttonationaleinkommen (BNE)} ist das Einkommen der \emph{Inländer},
    das sie durch Bereitstellung von Produktionsfaktoren im In- \emph{und} Ausland
    erzielt haben (\textbf{Inländerkonzept}) --- z.\,B. Pendler, Monteure, Künstler,
    Auslandsproduktion deutscher Unternehmen. „Inländer`` sind alle natürlichen Personen
    mit ständigem Wohnsitz im Inland (unabhängig von der Staatsbürgerschaft) und alle im
    Inland wirtschaftlich tätigen Unternehmen (auch in ausländischem Besitz).
\end{itemize}

\begin{merksatz}[title={\bfseries Formeln: Bruttowertschöpfung und BNE}]
\textbf{Bruttowertschöpfung} = (Brutto-)Produktionswerte $-$ Vorleistungen
\emph{(vermeidet Doppelzählungen!)}\\[0.4em]
\textbf{BIP} $-$ Inlandseinkommen von Nichtinländern $+$ Auslandseinkommen von Inländern
= \textbf{BNE}\\
(die beiden Korrekturposten bilden den \emph{Saldo der Primäreinkommen aus der übrigen
Welt})
\end{merksatz}

\begin{beispiel}
Ein Küchengerätehersteller bezieht Elektromotoren von einem Zulieferer. Würde man beim
BIP sowohl die Leistung des Motorenherstellers als auch die des Küchengeräteherstellers
voll zählen, wäre der Motor \emph{doppelt} erfasst --- denn im Preis der Rührmaschine
steckt der Motorpreis bereits. Deshalb werden von den Produktionswerten die
\textbf{Vorleistungen} abgezogen.
\end{beispiel}

Inlandsprodukt und Nationaleinkommen können \emph{brutto} (inkl. Abschreibungen) oder
\emph{netto} (nach Abzug der Abschreibungen) dargestellt werden.

\subsection{Nominales und reales BIP}

Das \textbf{nominale BIP} wird in laufenden/aktuellen Preisen (Marktpreisen) ausgewiesen.
Steigt es um 5~\%, kann das an gestiegenen \emph{Mengen}, an gestiegenen \emph{Preisen}
(Inflation) oder an beidem liegen. Um den \emph{echten} Leistungszuwachs zu ermitteln,
errechnet das Statistische Bundesamt das \textbf{reale BIP}: Es legt die Preise eines
\emph{festgelegten Basisjahres} zugrunde und eliminiert so die Preissteigerungen. Die
Wachstumsrate des realen BIP zeigt an, um wie viel Prozent die Produktion im Vergleich
zum Vorjahr gestiegen ist. Im internationalen Vergleich werden stets die Werte des
\emph{realen} BIP bzw. BNE zugrunde gelegt.

\subsection{Entstehung, Verwendung und Verteilung des BIP}

Das BIP kann über drei Rechnungen ermittelt werden --- alle führen zum gleichen Ergebnis:

\begin{center}
\begin{tabular}{p{3.4cm} p{4.6cm} p{6.4cm}}
\toprule
\textbf{Rechnung} & \textbf{Leitfrage} & \textbf{Schema (vereinfacht)}\\
\midrule
Entstehungs\-rechnung & \emph{Wo} ist die Wirtschaftsleistung entstanden? & Bruttowertschöpfung der Wirtschaftsbereiche $+$ Gütersteuern abzgl. Gütersubventionen $=$ BIP\\
Verwendungs\-rechnung & \emph{Wofür} wurde die Inlandsproduktion verwendet? & Private Konsumausgaben $+$ Konsumausgaben des Staates $+$ Bruttoinvestitionen $+$ Außenbeitrag (Exporte $-$ Importe) $=$ BIP\\
Verteilungs\-rechnung & \emph{An wen} fließen die Einkommen der Inländer? & Arbeitnehmerentgelt $+$ Unternehmens- und Vermögenseinkommen $=$ \textbf{Volkseinkommen}; $+$ Produktions-/Importabgaben abzgl. Subventionen $+$ Abschreibungen $-$ Saldo der Primäreinkommen aus der übrigen Welt $=$ BIP\\
\bottomrule
\end{tabular}
\end{center}

Hinweise zur Entstehungsrechnung: Leistungen der Unternehmen werden zu \emph{Marktpreisen}
bewertet; die (unverkaufte) Produktionsleistung des \emph{Staates} (Bildung, Polizei) zu
den \emph{Kosten der eingesetzten Produktionsfaktoren}. Bei privaten Haushalten zählen nur
Leistungen, für die Produktionskosten durch fremde Personen (z.\,B. Haushaltshilfen)
entstanden sind. Bei der Verteilungsrechnung wechselt die Bezugsgrundlage vom
\emph{Inland} zu den \emph{Inländern}.

\subsection{Grenzen der Aussagekraft des BIP}

\begin{itemize}
  \item \textbf{Eigenleistungen} privater Haushalte gelten nicht als Produktion;
    \textbf{ehrenamtliche} und andere unentgeltliche Tätigkeiten (z.\,B. Hausarbeit)
    bleiben unberücksichtigt.
  \item Die \textbf{Schwarzarbeit} (Teil der \emph{Schattenwirtschaft}) wird vom
    Statistischen Bundesamt nur \emph{geschätzt}.
  \item Erfasst werden auch Transaktionen, die den \emph{Wohlstand nicht erhöhen} (z.\,B.
    Reparatur von Unfall- und Umweltschäden zählt positiv zum BIP); umgekehrt bleiben
    wohlstandserhöhende, aber unbezahlte Leistungen unerfasst.
  \item Das BIP sagt nichts über die \emph{Verteilung} des Wohlstands, Umweltbelastungen
    oder Lebensqualität aus --- ein hohes Pro-Kopf-Einkommen schließt Armut von Teilen
    der Bevölkerung nicht aus.
\end{itemize}

\section{Primär- und Sekundärverteilung des Volkseinkommens}

\subsection{Primärverteilung: Lohn- und Gewinnquote}

Das über die Verteilungsrechnung errechnete \textbf{Volkseinkommen} ist die Summe aller
Arbeitnehmerentgelte sowie Unternehmens- und Vermögenseinkommen, die Inländern aus dem
In- und Ausland zugeflossen sind. Diese Verteilung heißt \textbf{Primärverteilung}, die
Einkommen heißen \emph{Primäreinkommen}. Sie basiert auf dem \textbf{Leistungsprinzip}:
Grundlage ist die wirtschaftliche Betätigung und die Bereitstellung von
Produktionsfaktoren. Man betrachtet die Verteilung \emph{funktional} (nach
Produktionsfaktoren: Lohnquote, Gewinnquote) oder \emph{personell} (nach Haushalten:
Pro-Kopf-Einkommen, Lorenzkurve).

\begin{merksatz}[title={\bfseries Formeln: Lohnquote und Gewinnquote}]
\[
\text{Lohnquote in \%} = \frac{\text{Arbeitnehmerentgelt}}{\text{Volkseinkommen}} \cdot 100
\qquad
\text{Gewinnquote in \%} = \frac{\text{Unternehmens- und Vermögenseinkommen}}{\text{Volkseinkommen}} \cdot 100
\]
2018 betrug die Lohnquote rund 69~\%, die Gewinnquote rund 31~\%. Beide Quoten ergänzen
sich stets zu 100~\%.
\end{merksatz}

Eine \emph{steigende Lohnquote} deutet an, dass der Anteil des Arbeitnehmereinkommens
zulasten der Unternehmens- und Vermögenseinkommen gestiegen ist. Vorsicht bei der
Interpretation: Auch Arbeitnehmer haben Zinseinkünfte (zählen zur Gewinnquote), und
Vorstandsgehälter zählen zum Arbeitnehmerentgelt (Lohnquote).

Die \textbf{personelle Einkommensverteilung} fragt, wie gleichmäßig das Volkseinkommen
auf die privaten Haushalte verteilt ist --- unabhängig von der Einkommensquelle.
Vereinfachter Verteilungsschlüssel ist das \textbf{Pro-Kopf-Einkommen} (Volkseinkommen
$\div$ Einwohnerzahl), das häufig als Wohlstandsindikator im Ländervergleich dient. Die
Verteilung wird mit der \textbf{Lorenzkurve} dargestellt: Sie zeigt, wie viel Prozent der
Einkommensempfänger wie viel Prozent des Volkseinkommens verdienen. Je weiter die Kurve
von der Diagonalen (Gleichverteilung) entfernt ist, desto ungleicher die Verteilung.

\subsection{Sekundärverteilung: Einkommensumverteilung}

Die Primärverteilung ist ungleich und wird häufig als ungerecht empfunden. Der Staat
korrigiert sie durch die \textbf{Sekundärverteilung} nach dem \textbf{Bedarfsprinzip}:
staatliche Umverteilung von Teilen des Primäreinkommens mit dem Ziel einer gerechteren
Einkommensverteilung. Ansatzpunkte:

\begin{itemize}
  \item \textbf{Progressive Einkommensteuer:} Wer über die Primärverteilung ein hohes
    Einkommen erzielt, wird zu einem höheren Prozentsatz besteuert.
  \item \textbf{Sozialversicherungsbeiträge} und \textbf{Sozialleistungen des Staates:}
    Wer keine Faktorleistung erbringen kann, erhält \emph{Transfereinkommen} --- z.\,B.
    Arbeitslosengeld, Renten, Wohngeld, Kindergeld.
\end{itemize}

\begin{merksatz}[title={\bfseries Formel: Verfügbares Einkommen der privaten Haushalte}]
\begin{tabular}{ll}
  & Faktoreinkommen der privaten Haushalte (Primäreinkommen)\\
$-$ & direkte Steuern\\
$-$ & Sozialversicherungsbeiträge\\
$+$ & Sozialleistungen des Staates (Transferzahlungen)\\
\midrule
$=$ & \textbf{verfügbares Einkommen} der privaten Haushalte
\end{tabular}\\[0.4em]
Das verfügbare Einkommen steht für \emph{Konsum} und \emph{Sparen} zur Verfügung ---
es ist das Ergebnis der Sekundärverteilung.
\end{merksatz}

Auch die \textbf{Gewerkschaften} beeinflussen die Einkommensverteilung über die
Tarifpolitik (Tariflöhne, Lohnsteigerungen über dem Produktivitätsfortschritt,
überproportionale Steigerungen für untere Lohngruppen). Zu hohe Forderungen können jedoch
die Stabilität der Unternehmen, ihre internationale Wettbewerbsfähigkeit und damit
Arbeitsplätze gefährden.

\section{Ziele der Stabilitätspolitik}

\subsection{Das Stabilitätsgesetz und das magische Viereck}

Das \textbf{Gesetz zur Förderung der Stabilität und des Wachstums der Wirtschaft (StabG)}
von 1967 gilt als \emph{Grundgesetz der Wirtschaftspolitik}. Nach \S\,1 StabG haben Bund
und Länder bei ihren wirtschafts- und finanzpolitischen Maßnahmen die Erfordernisse des
\textbf{gesamtwirtschaftlichen Gleichgewichts} zu beachten. Die Maßnahmen sollen im Rahmen
der marktwirtschaftlichen Ordnung \emph{gleichzeitig} beitragen zu:

\begin{enumerate}
  \item \textbf{Stabilität des Preisniveaus},
  \item \textbf{hohem Beschäftigungsstand},
  \item \textbf{außenwirtschaftlichem Gleichgewicht},
  \item \textbf{stetigem und angemessenem Wirtschaftswachstum}.
\end{enumerate}

Da es nahezu unmöglich ist, alle vier Ziele gleichzeitig vollständig zu erreichen, spricht
man vom \textbf{magischen Viereck}. Seine Ziele sind \emph{quantitativ}, d.\,h. der Grad
der Zielerreichung ist über Kennzahlen \emph{messbar}. Ergänzt um die \emph{qualitativen}
Ziele \textbf{Erhaltung einer lebenswerten Umwelt (Umweltschutz)} und \textbf{gerechte
Einkommens- und Vermögensverteilung} spricht man vom \textbf{magischen Sechseck} --- für
die qualitativen Ziele gibt es keinen klaren Maßstab der Zielerreichung.

Der \textbf{Jahreswirtschaftsbericht} (nach \S\,2 StabG zu Jahresbeginn dem Bundestag
vorzulegen) informiert über die wirtschaftlichen Zukunftserwartungen der Bundesregierung,
formuliert die wirtschafts- und finanzpolitischen Ziele für das laufende Jahr und nimmt
Stellung zum Jahresgutachten des \emph{Sachverständigenrates}.

\subsection{Die vier Ziele und ihre Messgrößen}

\begin{center}
\begin{tabular}{p{4.2cm} p{4.6cm} p{5.4cm}}
\toprule
\textbf{Ziel} & \textbf{Messgröße} & \textbf{Zielwert (Faustregel)}\\
\midrule
Hoher Beschäftigungsstand (Vollbeschäftigung) & Arbeitslosenquote (Bundesagentur für Arbeit) & erreicht bei ca. 3--4~\% (Sockelarbeitslosigkeit bleibt)\\
Preisniveaustabilität & Inflationsrate auf Basis des Verbraucherpreisindex (Warenkorb, Statistisches Bundesamt) & Inflationsrate knapp unter 2~\% (EZB-Definition, HVPI)\\
Außenwirtschaftliches Gleichgewicht & Außenbeitrag (Saldo aus Handels- und Dienstleistungsbilanz) & positiver Außenbeitrag von 1--2~\% des nominalen BIP\\
Stetiges und angemessenes Wirtschaftswachstum & jährliche Zuwachsrate des \emph{realen} BIP & Durchschnitt rund 2~\% pro Jahr, ohne heftige Konjunkturausschläge\\
\bottomrule
\end{tabular}
\end{center}

Erläuterungen: Die \emph{Arbeitslosenquote} setzt registrierte Arbeitslose ins Verhältnis
zu den zivilen Erwerbspersonen. \emph{Preisniveaustabilität} bedeutet nicht, dass kein
Preis steigen darf --- einzelne Preissteigerungen sollen durch Preissenkungen bei anderen
Gütern ausgeglichen werden. Der \emph{Außenbeitrag} (Exporte $-$ Importe von Waren und
Dienstleistungen) zeigt bei positivem Wert, dass die Volkswirtschaft mehr produziert als
verbraucht hat (Netto-Güterexport). \emph{Wirtschaftswachstum} wird stets an der
Zuwachsrate des realen BIP gemessen.

\begin{eselsbruecke}
Magisches Viereck: „\textbf{P}reise, \textbf{A}rbeit, \textbf{W}achstum,
\textbf{A}ußenhandel`` --- \textbf{PAWA}. Für das Sechseck kommen \textbf{U}mwelt und
\textbf{V}erteilung hinzu.
\end{eselsbruecke}

\subsection{Zielkonflikte und Zielharmonien}

Die Ziele stehen in unterschiedlichen Beziehungen zueinander --- deshalb „magisch``:

\begin{itemize}
  \item \textbf{Konkurrierende Zielbeziehungen (Zielkonflikte):}
    \emph{Wirtschaftswachstum vs. Preisniveaustabilität} (starkes Wachstum zieht i.\,d.\,R.
    Preiserhöhungen nach sich); \emph{Wirtschaftswachstum vs. Ökologie} (Wachstum um jeden
    Preis belastet die Umwelt); \emph{gerechte Verteilung vs. Wachstum} (zu starke
    Nivellierung der Einkommensunterschiede mindert Leistungsanreize und kann das
    Wachstum bremsen).
  \item \textbf{Komplementäre Zielbeziehungen (Zielharmonien):}
    \emph{Wirtschaftswachstum und Vollbeschäftigung} (steigt das BIP, werden mehr
    Produktionsfaktoren benötigt --- die Arbeitslosigkeit geht zurück);
    \emph{Wirtschaftswachstum und Umweltschutz} können harmonieren, wenn boomende
    Unternehmen in umweltfreundliche Produktionsverfahren investieren.
  \item \textbf{Indifferente Zielbeziehungen:} Ziele beeinflussen sich nicht ---
    z.\,B. außenwirtschaftliches Gleichgewicht und Erhaltung der Umwelt.
\end{itemize}

\begin{merksatz}[title={\bfseries Wichtig für die Prüfungsvorbereitung}]
Sie sollten sicher können: die \emph{Funktionen des Wettbewerbs}, die \emph{drei Säulen
des GWB} mit Kartellarten und Beispielen, Höchst- und Mindestpreise \emph{im
Preis-Mengen-Diagramm} einzeichnen und beurteilen, die \emph{Verwendungs- und
Verteilungsrechnung} des BIP schematisch darstellen, \emph{nominales und reales BIP}
unterscheiden, Kritikpunkte an der BIP-Berechnung aufzählen, die \emph{Lohnquote}
berechnen und das \emph{verfügbare Einkommen} ermitteln sowie die Ziele des magischen
Vierecks/Sechsecks nennen, ihre Messbarkeit beschreiben und Zielkonflikte bzw.
Zielharmonien mit Beispielen erläutern. Üben Sie mit dem Aufgabenheft und dem Quiz!
\end{merksatz}

\vspace{1em}
\hrule
\vspace{0.8em}
\noindent\textsf{\small Quellen: Rahmenplan Wirtschaftsbezogene Qualifikationen (DIHK), Abschnitte 1.1.1.2, 1.1.1.3, 1.1.2, 1.1.3.1;
Textband Volkswirtschaft (DIHK-Bildungs-GmbH), Kap.~1.2--1.3 (S.~12--18), Kap.~2 (S.~19--27), Kap.~3.1 (S.~29--38).
Nächste Lektion (Sa, 29.08.2026): Konjunktur, Wirtschaftspolitik \& Außenwirtschaft.}

\end{document}
